% Authors: Dr. Denys Dutykh (Mathematics Department, Khalifa University of Science and Technology,
% Abu Dhabi, UAE) and Prof. Laurent Vuillon (Univ. Savoie Mont Blanc, CNRS, LAMA, Chambery, France)
\section{Covering the unbounded tails}
\label{sec:tails}

A finite denominator check cannot establish weak minimality along an entire
Farey row. The missing step is a uniform separation between the source
ratio and every shorter inverse-Euclidean path. We obtain this separation
by covering an infinite lattice of gap lengths with finitely many boxes.
On each box the normalized ratio is controlled by its two values at a
matched pair of corners, and the reason two values suffice is an order
theorem for that ratio, Fibonacci denominator included.

\subsection{A common eigenchannel for every defect count}

Use the inverse-Euclidean matrices
$$
 \mathsf A=\begin{pmatrix}1&0\\1&1\end{pmatrix},\qquad
 \mathsf B=\begin{pmatrix}0&1\\1&1\end{pmatrix},\qquad z_0=\binom12.
$$
For $k\geqslant1$, a path with exactly $k$ defects is described uniquely
by nonnegative gaps $g=(g_0,\ldots,g_k)$: apply $\mathsf B^{g_i}$ and
then $\mathsf A$ for $i<k$, and apply only $\mathsf B^{g_k}$ at the end.
The order in this description is chronological. If $S=\sum_i g_i$, the
cost is $c=S+k+1$.

Put $\varphi=(1+\sqrt5)/2$, $\theta=\varphi^{-2}$, and $q=-\theta=\varphi-2$.
The matrix
$$
 T=\begin{pmatrix}1&\varphi\\\varphi&-1\end{pmatrix}
$$
satisfies $T^2=(1+\varphi^2)I$. Moreover
$$
 \varphi^{-g}T\mathsf B^gT^{-1}=\operatorname{diag}(1,q^g),
 \qquad
 \sqrt5\varphi^{-1}T\mathsf AT^{-1}
 =\begin{pmatrix}2&1\\q&-2q\end{pmatrix}.
$$
Thus we may start with $(X,Y)=(\varphi^2,2-\sqrt5)$, apply
\begin{equation}
 \binom XY\longmapsto
 \begin{pmatrix}2&q^{g_i}\\q&-2q^{g_i+1}\end{pmatrix}\binom XY
 \quad(i<k),
 \qquad (X,Y)\longmapsto(X,q^{g_k}Y)
 \quad(i=k),
 \label{eq:tail-channel}
\end{equation}
and denote the final $X^2+Y^2$ by $\Pi_k(g)$. The initial vector is
$\varphi^{-1}Tz_0$, so the product of all removed positive factors gives
$$
 \Pi_k(g)=5^k\varphi^{-2c}(1+\varphi^2)N(g).
$$
Binet's formula, with $2c+3$ odd, consequently gives
\begin{equation}
 \mathcal R_k(g):=\frac{N(g)}{F_{2c+3}}
 =\frac{\Pi_k(g)}{K_k(1-q^{2k+5+2S})},
 \qquad K_k=5^k\varphi^4.
 \label{eq:tail-normalization}
\end{equation}
The denominator is positive. This formula includes zero gaps and the
special final gap, which has no following defect matrix.

The right-hand side of \eqref{eq:tail-normalization} depends on the gaps
only through the numbers $x_j=q^{g_j}$ for $j<k$, which enter the matrices
\eqref{eq:tail-channel}, and $y=\theta^{2g_k}$, since $q^{2S}=y\prod_{j<k}x_j^2$.
It is a rational function of these variables whose denominator stays
positive when any of them is set to zero. We use it to give a meaning to
infinite gaps: the value at $g_j=\infty$ is the value at $x_j=0$, or at
$y=0$ for the final gap. An even gap has $x_j\in(0,1]$, an odd gap has
$x_j\in[-\theta,0)$, and the final gap has $y\in(0,1]$; with the infinite
values added, the ranges become the closed intervals $[0,1]$, $[-\theta,0]$
and $[0,1]$.

\subsection{Order in each parity coordinate}

\begin{theorem}[Normalized parity order]
\label{thm:tail-parity}
For every $k\geqslant1$, hold all gaps but $g_j$ fixed and increase $g_j$
by two. The normalized ratio $\mathcal R_k$ strictly increases if $j<k$
and $g_j$ is even. It strictly decreases if $j<k$ and $g_j$ is odd, or
if $j=k$ with either parity. More precisely, on the closed ranges just
described, $\mathcal R_k$ is a strictly decreasing function of each $x_j$,
$j<k$, and a strictly increasing function of $y$, whatever the values of the
other variables in their closed ranges; in particular the monotonicity
persists when some of the other gaps are infinite.
\end{theorem}

\begin{proof}
\emph{The state before a gap.} Just before any gap, the physical
column has second-to-first coordinate ratio $\xi\geqslant2$. This is true at
$(1,2)$; every intervening $\mathsf B$ block leaves a ratio greater than one,
and the following $\mathsf A$ adds one. Its channel therefore has $X>0$,
$Y<0$, and
$$
 h=-Y/X=\frac{\xi-\varphi}{1+\varphi\xi},\qquad
 \varphi^{-5}\leqslant h<\varphi^{-1}.
$$
This includes the first gap.

\emph{A nonfinal gap.} Write the two columns of the channel defect matrix as
$\mathsf u=(2,q)^{\mathsf T}$ and $\mathsf v=(1,-2q)^{\mathsf T}$. Their
physical preimages are
$$
 T^{-1}\mathsf u=\binom{\varphi^{-2}}1=:\mathsf a,
 \qquad T^{-1}\mathsf v=\binom{\varphi^{-1}}{\varphi^{-3}}=:\mathsf b,
 \qquad \varphi^{-3}\mathsf a\leqslant\mathsf b\leqslant\varphi\mathsf a.
$$
Let $S$ be the product of the physical matrices that follow the defect, a
matrix with nonnegative entries, and put $U=S\mathsf a$ and $V=S\mathsf b$.
Componentwise $\varphi^{-3}U\leqslant V\leqslant\varphi U$, all entries being
nonnegative, so that
$$
 \Delta:=\langle U,V\rangle\geqslant\varphi^{-3}\lVert U\rVert^2>0,
 \qquad \nu:=\lVert V\rVert^2\leqslant\varphi\Delta ;
$$
here $U\neq0$ because every factor of $S$ maps positive columns to positive
columns. The gap and the defect send the channel vector $(X,-hX)$ to
$X(\mathsf u-xh\mathsf v)$, with $x=x_j=q^{g_j}$, and the channel matrices that
follow are positive multiples of $TST^{-1}$. Since $T$ is scaled orthogonal,
$T^{\mathsf T}T=(1+\varphi^2)I$, the numerator is therefore a positive multiple,
independent of $x$, of
$$
 \Pi(x)=X^2\bigl(\lVert U\rVert^2-2xh\Delta+x^2h^2\nu\bigr).
$$

\emph{An even gap.} Here $0\leqslant x\leqslant1$, and $xh\nu<\Delta$ on this
whole closed range, because $h\nu<\varphi^{-1}\varphi\Delta=\Delta$. The
derivative $2X^2h(xh\nu-\Delta)$ of $\Pi$ is therefore negative, and $\Pi$
decreases strictly with $x$. The normalizing denominator
$1-q^{2k+5+2S}=1+\theta^{2k+5}\prod_ix_i^2\cdot y$ increases with $x^2$, so
$\mathcal R_k$ decreases strictly with $x$ on $[0,1]$. Increasing an even gap
by two decreases $x$, and so increases $\mathcal R_k$. This proves the even
case.

\emph{An odd gap.} Here $x=-\tau$ with $0\leqslant\tau\leqslant\theta$. Set
$\zeta=\theta^{2k+5+2\sum_{i\ne j}g_i}$, so that the normalizing denominator is
$1+\zeta\tau^2$. Up to the fixed positive factor $K_k$, the normalized
function is
$$
 F(\tau)=\frac{\omega_0+2\omega_1\tau+\omega_2\tau^2}{1+\zeta\tau^2},\qquad
 \omega_0=X^2\lVert U\rVert^2,\quad \omega_1=X^2h\Delta,\quad \omega_2=X^2h^2\nu.
$$
From the bounds on $\Delta$ and $h$, $\omega_0/\omega_1\leqslant\varphi^8=\theta^{-4}$.
Differentiation gives
$$
 F'(\tau)=\frac{2(\omega_1+\omega_2\tau-\omega_0\zeta\tau-\omega_1\zeta\tau^2)}{(1+\zeta\tau^2)^2},
$$
and the numerator is positive on the closed range: using $\omega_2\tau\geqslant0$,
$\tau\leqslant\theta$ and $\zeta\leqslant\theta^{2k+5}$,
$$
 \omega_1+\omega_2\tau-\omega_0\zeta\tau-\omega_1\zeta\tau^2
 \geqslant \omega_1\bigl(1-\theta^{-4}\theta^{2k+5}\theta-\theta^{2k+5}\theta^2\bigr)
 =\omega_1\bigl(1-\theta^{2k+2}-\theta^{2k+7}\bigr)>0.
$$
Indeed $\theta<1/2$ and $k\geqslant1$, so even the boundary case has
$\theta^4+\theta^9<1$. Thus $F$ increases strictly with $\tau$, that is,
$\mathcal R_k$ decreases strictly with $x=-\tau$ on $[-\theta,0]$. Increasing the
odd gap by two decreases $\tau$, and therefore decreases $\mathcal R_k$.

\emph{The final gap.} Put $\zeta=\theta^{2k+5+2\sum_{i<k}g_i}$. Up to $K_k$, the
normalized function of $y=\theta^{2g_k}$ is
$$
 \frac{X^2+Y^2y}{1+\zeta y},
$$
whose derivative in $y$ is $(Y^2-X^2\zeta)/(1+\zeta y)^2$, of the sign of
$h^2-\zeta$. This is positive since $h^2\geqslant\theta^5>\theta^{2k+5}\geqslant\zeta$,
including $k=1$. So $\mathcal R_k$ increases strictly with $y$ on $[0,1]$;
increasing the final gap decreases $y$ and proves the last assertion.

\emph{The other variables anywhere in their ranges.} The three derivative
computations used only the bounds on $h$, $\Delta$, $\nu$ and $\zeta$, and these
hold whatever the values of the other variables in their closed ranges, not
only at the values $q^{g}$, $\theta^{2g}$ and $0$. Indeed, for
$-\theta\leqslant x\leqslant1$ the channel action $\operatorname{diag}(1,x)$ of a gap
corresponds in physical coordinates to
$$
 M_x=T^{-1}\operatorname{diag}(1,x)\,T=\frac1{1+\varphi^2}
 \begin{pmatrix}1+\varphi^2x&\varphi(1-x)\\\varphi(1-x)&\varphi^2+x\end{pmatrix},
$$
whose entries are nonnegative because $x\geqslant-\theta=-\varphi^{-2}$ and
$x\leqslant1$; it is $\varphi^{-g}\mathsf B^g$ when $x=q^g$, and for $x=0$ it is the
limit that represents an infinite gap. It maps positive columns to positive
columns. The final gap enters only through $y$, and acts on the channel as
$\operatorname{diag}(1,\sqrt y)$ with $0\leqslant\sqrt y\leqslant1$, a matrix of the same
kind. Every suffix is thus still a product of nonnegative matrices, and the
bounds on $\Delta$ and $\nu$ hold as shown above. Next, applied to a column
$(1,\xi)^{\mathsf T}$ with $\xi\geqslant2$, the matrix $M_x$ gives a column with
positive entries whose second coordinate exceeds the first by
$$
 \frac{\varphi^{-1}+\xi+x\varphi^2(\xi-\varphi)}{1+\varphi^2}
 \geqslant\frac{\varphi^{-1}+\varphi}{1+\varphi^2}>0 ,
$$
since for $x<0$ one has $x\varphi^2(\xi-\varphi)\geqslant-\theta\varphi^2(\xi-\varphi)=-(\xi-\varphi)$;
the next $\mathsf A$ therefore gives a ratio above $2$ again, and the bounds on
$h$ hold before every gap. Finally, every other variable has absolute value at
most $1$, so $\zeta\leqslant\theta^{2k+5}$ in both places where it was used. Each of
the three derivative computations therefore applies at every point of the
closed ranges, and the monotonicity is strict there. With $x_i=0$ or $y=0$
this is the case of infinite gaps, which the covers below use.
\end{proof}

\Cref{fig:parity-order} shows the theorem at work in the simplest stratum,
one defect, where the ratio depends on two gaps only: along a nonfinal gap the
even and the odd values form two monotone sequences that approach the common
limit from opposite sides, and along the final gap the ratio decreases.

\begin{figure}[tb]
\centering
% Authors: Dr. Denys Dutykh (Mathematics Department, Khalifa University of Science and
% Technology, Abu Dhabi, UAE) and Prof. Laurent Vuillon (Univ. Savoie Mont Blanc, CNRS, LAMA)
% Generated by supplement/code/make_figures.py from exact data; do not edit by hand.
\begin{tikzpicture}[font=\footnotesize, >={Stealth[length=4pt]}]
\node[ink, right] at (-0.1,4.45) {(a) a nonfinal gap};
\draw[ink, line width=0.5pt] (0,0) -- (5.2,0);
\draw[ink, line width=0.5pt] (0,0) -- (0,4);
\draw[ink, line width=0.5pt] (0.236,0) -- (0.236,-0.08);
\node[ink, below] at (0.236,-0.08) {$0$};
\draw[ink, line width=0.5pt] (1.024,0) -- (1.024,-0.08);
\node[ink, below] at (1.024,-0.08) {$2$};
\draw[ink, line width=0.5pt] (1.812,0) -- (1.812,-0.08);
\node[ink, below] at (1.812,-0.08) {$4$};
\draw[ink, line width=0.5pt] (2.6,0) -- (2.6,-0.08);
\node[ink, below] at (2.6,-0.08) {$6$};
\draw[ink, line width=0.5pt] (3.388,0) -- (3.388,-0.08);
\node[ink, below] at (3.388,-0.08) {$8$};
\draw[ink, line width=0.5pt] (4.176,0) -- (4.176,-0.08);
\node[ink, below] at (4.176,-0.08) {$10$};
\draw[ink, line width=0.5pt] (4.964,0) -- (4.964,-0.08);
\node[ink, below] at (4.964,-0.08) {$12$};
\draw[ink, line width=0.5pt] (0,0) -- (-0.08,0);
\node[ink, left] at (-0.08,0) {$10^{-7}$};
\draw[ink, line width=0.5pt] (0,1.231) -- (-0.08,1.231);
\node[ink, left] at (-0.08,1.231) {$10^{-5}$};
\draw[ink, line width=0.5pt] (0,2.462) -- (-0.08,2.462);
\node[ink, left] at (-0.08,2.462) {$10^{-3}$};
\draw[ink, line width=0.5pt] (0,3.692) -- (-0.08,3.692);
\node[ink, left] at (-0.08,3.692) {$10^{-1}$};
\node[ink] at (2.6,-0.75) {nonfinal gap $g_0$};
\node[ink, rotate=90] at (-1.15,2) {distance to the limit};
\draw[slate, line width=0.5pt] (0.236,3.556) -- (0.63,3.308) -- (1.024,3.047) -- (1.418,2.792) -- (1.812,2.534) -- (2.206,2.277) -- (2.6,2.02) -- (2.994,1.762) -- (3.388,1.505) -- (3.782,1.248) -- (4.176,0.991) -- (4.57,0.733) -- (4.964,0.476);
\filldraw[navy, fill=white, line width=0.7pt] (0.236,3.556) circle (1.5pt);
\filldraw[navy, fill=white, line width=0.7pt] (1.024,3.047) circle (1.5pt);
\filldraw[navy, fill=white, line width=0.7pt] (1.812,2.534) circle (1.5pt);
\filldraw[navy, fill=white, line width=0.7pt] (2.6,2.02) circle (1.5pt);
\filldraw[navy, fill=white, line width=0.7pt] (3.388,1.505) circle (1.5pt);
\filldraw[navy, fill=white, line width=0.7pt] (4.176,0.991) circle (1.5pt);
\filldraw[navy, fill=white, line width=0.7pt] (4.964,0.476) circle (1.5pt);
\filldraw[wine, fill=white, line width=0.7pt] (0.58,3.258) rectangle (0.68,3.358);
\filldraw[wine, fill=white, line width=0.7pt] (1.368,2.742) rectangle (1.468,2.842);
\filldraw[wine, fill=white, line width=0.7pt] (2.156,2.227) rectangle (2.256,2.327);
\filldraw[wine, fill=white, line width=0.7pt] (2.944,1.712) rectangle (3.044,1.812);
\filldraw[wine, fill=white, line width=0.7pt] (3.732,1.198) rectangle (3.832,1.298);
\filldraw[wine, fill=white, line width=0.7pt] (4.52,0.683) rectangle (4.62,0.783);
\filldraw[navy, fill=white, line width=0.7pt] (2.248,3.75) circle (1.5pt);
\node[navy, right, font=\scriptsize] at (2.498,3.75) {even $g_0$: below the limit};
\filldraw[wine, fill=white, line width=0.7pt] (2.198,3.3) rectangle (2.298,3.4);
\node[wine, right, font=\scriptsize] at (2.498,3.35) {odd $g_0$: above the limit};
\node[ink, right] at (6.85,4.45) {(b) the final gap};
\draw[ink, line width=0.5pt] (6.95,0) -- (12.15,0);
\draw[ink, line width=0.5pt] (6.95,0) -- (6.95,4);
\draw[ink, line width=0.5pt] (7.186,0) -- (7.186,-0.08);
\node[ink, below] at (7.186,-0.08) {$0$};
\draw[ink, line width=0.5pt] (7.974,0) -- (7.974,-0.08);
\node[ink, below] at (7.974,-0.08) {$2$};
\draw[ink, line width=0.5pt] (8.762,0) -- (8.762,-0.08);
\node[ink, below] at (8.762,-0.08) {$4$};
\draw[ink, line width=0.5pt] (9.55,0) -- (9.55,-0.08);
\node[ink, below] at (9.55,-0.08) {$6$};
\draw[ink, line width=0.5pt] (10.338,0) -- (10.338,-0.08);
\node[ink, below] at (10.338,-0.08) {$8$};
\draw[ink, line width=0.5pt] (11.126,0) -- (11.126,-0.08);
\node[ink, below] at (11.126,-0.08) {$10$};
\draw[ink, line width=0.5pt] (11.914,0) -- (11.914,-0.08);
\node[ink, below] at (11.914,-0.08) {$12$};
\draw[ink, line width=0.5pt] (6.95,0.315) -- (6.87,0.315);
\node[ink, left] at (6.87,0.315) {$0.79$};
\draw[ink, line width=0.5pt] (6.95,1.115) -- (6.87,1.115);
\node[ink, left] at (6.87,1.115) {$0.80$};
\draw[ink, line width=0.5pt] (6.95,1.915) -- (6.87,1.915);
\node[ink, left] at (6.87,1.915) {$0.81$};
\draw[ink, line width=0.5pt] (6.95,2.715) -- (6.87,2.715);
\node[ink, left] at (6.87,2.715) {$0.82$};
\draw[ink, line width=0.5pt] (6.95,3.515) -- (6.87,3.515);
\node[ink, left] at (6.87,3.515) {$0.83$};
\node[ink] at (9.55,-0.75) {final gap $g_1$};
\draw[navy, line width=0.8pt, ] (7.186,2.733) -- (7.58,0.634) -- (7.974,0.328) -- (8.368,0.283) -- (8.762,0.277) -- (9.156,0.276) -- (9.55,0.276) -- (9.944,0.276) -- (10.338,0.276) -- (10.732,0.276) -- (11.126,0.276) -- (11.52,0.276) -- (11.914,0.276);
\draw[wine, line width=0.8pt, densely dashed] (7.186,3.724) -- (7.58,1.771) -- (7.974,1.486) -- (8.368,1.444) -- (8.762,1.438) -- (9.156,1.437) -- (9.55,1.437) -- (9.944,1.437) -- (10.338,1.437) -- (10.732,1.437) -- (11.126,1.437) -- (11.52,1.437) -- (11.914,1.437);
\filldraw[navy, fill=white, line width=0.6pt] (7.186,2.733) circle (1.25pt);
\filldraw[navy, fill=white, line width=0.6pt] (7.58,0.634) circle (1.25pt);
\filldraw[navy, fill=white, line width=0.6pt] (7.974,0.328) circle (1.25pt);
\filldraw[navy, fill=white, line width=0.6pt] (8.368,0.283) circle (1.25pt);
\filldraw[navy, fill=white, line width=0.6pt] (8.762,0.277) circle (1.25pt);
\filldraw[navy, fill=white, line width=0.6pt] (9.156,0.276) circle (1.25pt);
\filldraw[navy, fill=white, line width=0.6pt] (9.55,0.276) circle (1.25pt);
\filldraw[navy, fill=white, line width=0.6pt] (9.944,0.276) circle (1.25pt);
\filldraw[navy, fill=white, line width=0.6pt] (10.338,0.276) circle (1.25pt);
\filldraw[navy, fill=white, line width=0.6pt] (10.732,0.276) circle (1.25pt);
\filldraw[navy, fill=white, line width=0.6pt] (11.126,0.276) circle (1.25pt);
\filldraw[navy, fill=white, line width=0.6pt] (11.52,0.276) circle (1.25pt);
\filldraw[navy, fill=white, line width=0.6pt] (11.914,0.276) circle (1.25pt);
\filldraw[wine, fill=white, line width=0.6pt] (7.186,3.724) circle (1.25pt);
\filldraw[wine, fill=white, line width=0.6pt] (7.58,1.771) circle (1.25pt);
\filldraw[wine, fill=white, line width=0.6pt] (7.974,1.486) circle (1.25pt);
\filldraw[wine, fill=white, line width=0.6pt] (8.368,1.444) circle (1.25pt);
\filldraw[wine, fill=white, line width=0.6pt] (8.762,1.438) circle (1.25pt);
\filldraw[wine, fill=white, line width=0.6pt] (9.156,1.437) circle (1.25pt);
\filldraw[wine, fill=white, line width=0.6pt] (9.55,1.437) circle (1.25pt);
\filldraw[wine, fill=white, line width=0.6pt] (9.944,1.437) circle (1.25pt);
\filldraw[wine, fill=white, line width=0.6pt] (10.338,1.437) circle (1.25pt);
\filldraw[wine, fill=white, line width=0.6pt] (10.732,1.437) circle (1.25pt);
\filldraw[wine, fill=white, line width=0.6pt] (11.126,1.437) circle (1.25pt);
\filldraw[wine, fill=white, line width=0.6pt] (11.52,1.437) circle (1.25pt);
\filldraw[wine, fill=white, line width=0.6pt] (11.914,1.437) circle (1.25pt);
\draw[navy, line width=0.8pt, ] (9.55,3.7) -- (9.95,3.7);
\filldraw[navy, fill=white, line width=0.6pt] (9.75,3.7) circle (1.25pt);
\node[navy, right, font=\scriptsize] at (10,3.7) {$g_0=2$};
\draw[wine, line width=0.8pt, densely dashed] (9.55,3.3) -- (9.95,3.3);
\filldraw[wine, fill=white, line width=0.6pt] (9.75,3.3) circle (1.25pt);
\node[wine, right, font=\scriptsize] at (10,3.3) {$g_0=3$};
\node[ink, rotate=90] at (5.8,2) {$N/F_{2c+3}$};
\end{tikzpicture}

\caption{The normalized parity order for one defect, computed exactly. (a)
The one-defect paths with final gap $0$, as a function of the gap $g_0$ before
the defect: the distance of the ratio $N/F_{2c+3}$ from its value
$3/(\sqrt5\,\varphi)=(15-3\sqrt5)/10$ at $g_0=\infty$, on a logarithmic axis.
The even values (navy circles) lie below that limit and increase towards it,
the odd values (wine squares) lie above it and decrease. (b) The ratio itself
as a function of the final gap $g_1$, for $g_0=2$ and $g_0=3$: it decreases
for either parity.}
\label{fig:parity-order}
\end{figure}

\subsection{A finite cover of the infinite gap lattice}

Fix the parity of every gap. A coordinate domain is either a finite
step-two interval or an infinite step-two tail. For a product box, its
lower matched corner uses the lower endpoint in increasing coordinates
and the upper endpoint in decreasing coordinates; the upper matched
corner makes the opposite choices. An infinite endpoint means $x_j=0$, or
$y=0$ for the final gap. The numerator and denominator in
\eqref{eq:tail-normalization} are always evaluated at the same corner.
Moving one coordinate at a time from a point of the box to a corner, and
applying \cref{thm:tail-parity} on the closed ranges at each move, gives
\begin{equation}
 \mathcal R_k(g^{\min})\leqslant\mathcal R_k(g)
 \leqslant\mathcal R_k(g^{\max}).
 \label{eq:tail-corner-bounds}
\end{equation}
A corner need not satisfy a subsequently imposed lower bound on the total
gap sum; it still safely bounds that restricted part of the box.

For a required total $S\geqslant S_0$, set
$\eta=\lceil S_0/(k+1)\rceil$. Every eligible tuple has some gap at least
$\eta$. Partition it by the first such gap: earlier coordinates are below
$\eta$, that coordinate is at least $\eta$, and later coordinates are
unrestricted. Intersecting with a fixed parity gives $k+1$ disjoint root
boxes whenever $\eta\geqslant2$. Empty parity intersections are discarded.
Their union contains every eligible tuple exactly once.

Every nonterminal box is split into two disjoint step-two domains along
one coordinate, namely the non-singleton coordinate whose domain has the
smallest lower endpoint, the first such coordinate in case of a tie: a finite
domain of $n$ points is cut into its first $\lfloor n/2\rfloor$ points and the
rest, and an infinite tail into its first point and the rest.
A finite box may be discarded if its maximum total is
less than $S_0$. Otherwise it closes only when its upper matched corner
is strictly below the lower source threshold, or its lower matched corner
is strictly above the upper source threshold. A singleton inside the
closed source band would be an explicit failure. A subdivision tree all of
whose leaves close therefore settles the whole infinite domain: the root
boxes cover the eligible tuples, each exactly once, each split partitions a
box, and every leaf lies entirely on one side of the band by
\eqref{eq:tail-corner-bounds}.
Whether a given tree closes is a finite question, settled for the trees of
this paper in \cref{prop:finite-verifications}.

All the quantities compared here lie in $\Z[\varphi]$: the channel
entries, every power of $q$, and $K_k$. Writing $a+b\varphi$ for the pair
of integers $(a,b)$ and clearing threshold denominators beforehand, the
sign of $a+b\varphi$ is the sign of $2a+b+b\sqrt5$, which is decided by
the signs of the two integers $2a+b$ and $b$ and, when these oppose each
other, by the sign of $(2a+b)^2-5b^2$. Every comparison below is therefore
an exact comparison of integers; no decimal approximation of $\sqrt5$
enters.

\subsection{The six unbounded source windows}

With the constant $C_P=\varphi^2\lambda^{P-1}$ of \cref{sec:row-labels}, which
the source labels of row $P$ approach, define the rational ceilings and
starting costs by
\begin{equation}
 \begin{array}{c@{\qquad}c@{\qquad}c@{\qquad}c}
 P&c_P&U_P&\text{middle defects}\\[2pt]
 8&34&707/1000&2\\
 9&41&2929/5000&2,3\\
 10&46&4857/10000&2,3,4\\
 11&84&20111/50000&2,3,4,5\\
 12&110&333511/10^6&2,3,4,5,6\\
 13&156&2765404267/10^{10}&2,3,4,5,6,7.
 \end{array}
 \label{eq:tail-source-windows}
\end{equation}
For row $P$, the closed source band is $[C_P,U_P]$. The paths with no
defect and with a single defect are excluded directly in \cref{sec:assembly},
by the Fibonacci maximum and by \cref{lem:one-defect-minimum}; the paths with
at least two defects are the business of the covers. For each pair $(P,k)$
with $2\leqslant k\leqslant P-6$ (the middle defects of the table) the cover
is built with $S_0=c_P-k-1$ and both thresholds; for $k=P-5$ with the lower
threshold alone. \Cref{thm:tail-covers} in \cref{sec:assembly} states what
the closed trees prove.

The covers prove only the source-band exclusions needed at the actual tail
junctions. They do not classify every extremal ray or assert the same bands
at smaller costs. That narrower target makes one uniform order theorem
sufficient for every stratum.
